\documentclass[11pt]{article}
\usepackage[margin=1in]{geometry}
\usepackage{booktabs}
\usepackage{longtable}
\usepackage{graphicx}
\usepackage{amsmath,amssymb}
\usepackage{natbib}
\usepackage{hyperref}
\usepackage[T1]{fontenc}

\title{Supplementary Material\\[4pt]\large Blockchain Access Control for Agricultural Edge IoT: A 61-Day Hyperledger Fabric Field Study and Implementation Postmortem (V10)}
\author{}
\date{}

% AUTO-GENERATED by analysis/scripts/11_build_macros.py -- DO NOT EDIT BY HAND.
% Every value here is computed from paper_iot/V10/analysis/data/raw by the
% numbered scripts in analysis/scripts/. Re-run analysis/scripts/run_all.sh
% to regenerate.
\newcommand{\DeploymentDays}{61}
\newcommand{\Sensors}{50}
\newcommand{\Gateways}{4}
\newcommand{\Zones}{4}
\newcommand{\AuthDecisions}{209,000}
\newcommand{\SensorReadings}{146,400}
\newcommand{\OrdinaryDenials}{4,860}
\newcommand{\OrdinaryDenialPct}{2.33}
\newcommand{\DecisionPathUptime}{99.36}
\newcommand{\GatewayHourUptime}{99.84}
\newcommand{\OutageCount}{3}
\newcommand{\OutageTotalHours}{9.3}
\newcommand{\OutageMinHours}{0.2}
\newcommand{\OutageMaxHours}{8.7}
\newcommand{\CheckAccessMeanMs}{280}
\newcommand{\CheckAccessSdMs}{40}
\newcommand{\SensorWriteMeanMs}{1,220}
\newcommand{\SensorWriteSdMs}{60}
\newcommand{\SensorWritePNinetyFiveMs}{1,319}
\newcommand{\SensorWritePNinetyNineMs}{1,360}
\newcommand{\LatencyGapMs}{940}
\newcommand{\RoleOpCells}{36}
\newcommand{\GrantRateMin}{90.6}
\newcommand{\GrantRateMax}{100.0}
\newcommand{\FieldRevocationDenials}{973}
\newcommand{\FieldRevocationPerDay}{17.1}
\newcommand{\PeerRunsTotal}{32}
\newcommand{\PeerRunsPerConfig}{8}
\newcommand{\PeerTxPerRun}{130}
\newcommand{\TotalLatencyPeerFour}{1,203.9}
\newcommand{\SpanLatencyPeerFour}{321.3}
\newcommand{\RemainingLatencyPeerFour}{882.6}
\newcommand{\PNinetyFiveLatencyPeerFour}{1,294.2}
\newcommand{\SpanSharePeerFour}{26.7}
\newcommand{\TotalLatencyPeerEight}{1,149.8}
\newcommand{\SpanLatencyPeerEight}{320.4}
\newcommand{\RemainingLatencyPeerEight}{829.4}
\newcommand{\PNinetyFiveLatencyPeerEight}{1,239.5}
\newcommand{\SpanSharePeerEight}{27.9}
\newcommand{\TotalLatencyPeerSixteen}{1,057.7}
\newcommand{\SpanLatencyPeerSixteen}{322.8}
\newcommand{\RemainingLatencyPeerSixteen}{734.9}
\newcommand{\PNinetyFiveLatencyPeerSixteen}{1,148.3}
\newcommand{\SpanSharePeerSixteen}{30.5}
\newcommand{\TotalLatencyPeerThirtyTwo}{812.8}
\newcommand{\SpanLatencyPeerThirtyTwo}{317.0}
\newcommand{\RemainingLatencyPeerThirtyTwo}{495.8}
\newcommand{\PNinetyFiveLatencyPeerThirtyTwo}{898.8}
\newcommand{\SpanSharePeerThirtyTwo}{39.0}
\newcommand{\TotalLatencyDiffFourVThirtyTwo}{391.0}
\newcommand{\TotalLatencyDiffFourVThirtyTwoCILow}{388.3}
\newcommand{\TotalLatencyDiffFourVThirtyTwoCIHigh}{393.8}
\newcommand{\RemainingLatencyDiffFourVThirtyTwo}{386.7}
\newcommand{\SpanDiffFourVThirtyTwo}{4.31}
\newcommand{\SpanDiffFourVThirtyTwoCINinetyLow}{4.03}
\newcommand{\SpanDiffFourVThirtyTwoCINinetyHigh}{4.60}
\newcommand{\SpanDiffFourVThirtyTwoWelchT}{28.77}
\newcommand{\SpanDiffFourVThirtyTwoP}{$<$0.001}
\newcommand{\SpanTostP}{$<$0.001}
\newcommand{\PNinetyFiveDiffFourVThirtyTwo}{395.4}
\newcommand{\LogLinearSlope}{-126.5}
\newcommand{\LogLinearRSquared}{0.891}
\newcommand{\LogLinearP}{0.056}
\newcommand{\SpanTrendSlope}{-1.05}
\newcommand{\SpanTrendP}{0.451}
\newcommand{\SpanTrendRSquared}{0.301}
\newcommand{\SeqEffectPSpanNFour}{0.735}
\newcommand{\SeqEffectPTotalNThirtyTwo}{0.923}
\newcommand{\PeerEffectPTotalLatency}{$<$0.001}
\newcommand{\BlockEffectPTotalLatency}{0.537}
\newcommand{\SpanLogTwoPeersJointP}{0.0017}
\newcommand{\ThroughputRunsTotal}{100}
\newcommand{\ThroughputConditionDiff}{5.65}
\newcommand{\ThroughputConditionCILow}{2.88}
\newcommand{\ThroughputConditionCIHigh}{8.41}
\newcommand{\ThroughputConditionP}{$<$0.001}
\newcommand{\ThroughputInteractionCoef}{-0.0021}
\newcommand{\ThroughputInteractionP}{0.97}
\newcommand{\ThroughputPositionP}{0.96}
\newcommand{\ThroughputModelRSquared}{0.317}
\newcommand{\ThroughputPctDiffMean}{9.4}
\newcommand{\ThroughputPctDiffSd}{1.2}
\newcommand{\ThroughputPctDiffMin}{6.9}
\newcommand{\ThroughputPctDiffMax}{10.7}
\newcommand{\ThroughputSigAfterHolm}{10}
\newcommand{\ThroughputBaselinePeak}{70.5}
\newcommand{\ThroughputHrbacPeak}{63.0}
\newcommand{\ThroughputPeakLevel}{40}
\newcommand{\ThroughputHrbacDeclinePct}{14.1}
\newcommand{\ThroughputCampaignSeqCorr}{0.995}
\newcommand{\SecurityAttemptsTotal}{8,000}
\newcommand{\SecurityDenialMechanisms}{6}
\newcommand{\OracleScoredAttempts}{2,666}
\newcommand{\OracleConfirmed}{852}
\newcommand{\OracleConfirmedPct}{32.0}
\newcommand{\OracleDiscrepant}{1,814}
\newcommand{\OracleDiscrepantPct}{68.0}
\newcommand{\CrtModuli}{97, 101, 103}
\newcommand{\CrtMinBound}{9,797}
\newcommand{\CrtExceedPct}{100.0000}
\newcommand{\CrtTwoResidueCount}{131,795}
\newcommand{\CrtTwoResiduePct}{90.0}
\newcommand{\CrtCorrectedModuli}{253, 254, 255}
\newcommand{\CrtCorrectedBound}{64,262}
\newcommand{\CrtCorrectedLegalMax}{51,455}
\newcommand{\QuantMaxErrorPct}{0.37}
\newcommand{\QuantRmsErrorPct}{0.21}
\newcommand{\SignatureValidPct}{100.0}
\newcommand{\RevocationStageRecords}{200}
\newcommand{\RevocationTraceSpanDays}{55}
\newcommand{\RevocationDelayedCount}{57}
\newcommand{\RevocationDelayedPct}{28.5}
\newcommand{\RevocationSubjects}{132}
\newcommand{\RevocationSubjectsAllStages}{0}
\newcommand{\RevocationSubjectsOnce}{87}
\newcommand{\EnergySingleChannelMj}{24.13}
\newcommand{\EnergyCrtMj}{12.42}
\newcommand{\EnergyReductionPct}{48.5}
\newcommand{\EnergyReductionCILow}{47.7}
\newcommand{\EnergyReductionCIHigh}{49.3}
\newcommand{\ShaTwoFiveSixMeanUs}{15.1}
\newcommand{\SignMeanUs}{894.8}
\newcommand{\VerifyMeanUs}{101.9}
\newcommand{\SignPipelineUs}{910}
\newcommand{\ReleasedRowsTotal}{370,118}
\newcommand{\ReleasedTracesCount}{12}


\begin{document}
\maketitle

This supplement accompanies the V10 manuscript. Every table here is generated by \texttt{paper\_iot/V10/analysis/scripts/} from the same raw-data package as the main text (\texttt{run\_all.sh} regenerates both).

\section*{S1. System-comparison summary}

Table~\ref{tab:s1} summarises, at a coarse level, how the systems discussed in Section~2 (Related work) report their evaluation, to support the claim that laboratory or simulated short-duration testbeds dominate the literature. This table is qualitative (drawn from each cited paper's own reported methodology) and is not generated from the raw-data package.

\begin{table}[h]
\centering
\caption{Evaluation methodology of related Fabric-based / blockchain access-control systems.}
\label{tab:s1}
\small
\begin{tabular}{p{3.4cm}p{5.0cm}p{4.0cm}}
\toprule
System & Evaluation setting & Duration \\
\midrule
Idrissi \& Palmieri~\citep{idrissi2024blockchain} & Constrained ARM testbed & Single-session bench \\
Tzenetopoulos et al.~\citep{tzenetopoulos2022hlfkubed} & Edge cluster (Raspberry Pi) & Single-session bench \\
Pajooh et al.~\citep{pajooh2021hyperledger} & Lab testbed & Single-session bench \\
Ke \& Park~\citep{ke2023performance} & Simulation & Single-session bench \\
Zaidi et al.~\citep{zaidi2023fabrication} & Lab testbed & Single-session bench \\
Bandara et al.~\citep{bandara2021fabric} & Lab testbed & Single-session bench \\
This work & Working farm, 61-day field deployment plus separate controlled campaigns & 61 days field + multi-day controlled campaigns \\
\bottomrule
\end{tabular}
\end{table}

\section*{S2. Full effective-permission matrix (deployed hierarchy)}

Generated directly from \texttt{chaincode/hrbac/roles.go} (the same computation used in \texttt{analysis/scripts/06\_security\_oracle.py}), not transcribed.

\begin{center}
\begin{tabular}{lcccccccccc}
\toprule
Role & \rotatebox{60}{ReadOwn} & \rotatebox{60}{ReadZone} & \rotatebox{60}{ReadAll} & \rotatebox{60}{WriteSensor} & \rotatebox{60}{ControlZone} & \rotatebox{60}{ControlAll} & \rotatebox{60}{ReadAudit} & \rotatebox{60}{ManageRoles} & \rotatebox{60}{IssueCrossZone} & \rotatebox{60}{AdminAll} \\
\midrule
Admin & Yes & Yes & Yes & Yes & Yes & Yes & Yes & Yes & Yes & Yes \\
Gateway &  & Yes &  & Yes & Yes &  &  &  &  &  \\
Farmer & Yes & Yes &  & Yes & Yes &  & Yes &  & Yes &  \\
Agronomist & Yes & Yes &  &  & Yes &  & Yes &  & Yes &  \\
Certifier & Yes & Yes &  &  &  &  & Yes &  &  &  \\
SupplyChain & Yes &  &  &  &  &  &  &  &  &  \\
Sensor &  &  &  & Yes &  &  &  &  &  &  \\
\bottomrule
\end{tabular}
\end{center}

\section*{S3. Peer-scaling topology parameters}

All four configurations (4, 8, 16 and 32 logical peers) run on the same four physical Raspberry Pi hosts identified by a single \texttt{host\_id} value in \texttt{raw/experiments/peer\_scaling\_runs.csv}. Endorsement policy, orderer topology, batch parameters and CPU/memory allocation are held constant across configurations per the deployment configuration in Table~2 of the main text; the released run manifest does not record per-run CPU, memory, disk, block height or CouchDB document-count telemetry, so pre-run system state cannot be characterised beyond what Table~S7 below reports.

\section*{S4. Experimental unit by comparison}

\begin{table}[h]
\centering
\caption{Experimental unit used for every controlled-campaign comparison in this article.}
\label{tab:s4}
\small
\begin{tabular}{p{5.2cm}p{4.6cm}p{3.0cm}}
\toprule
Comparison & Unit & $n$ \\
\midrule
Peer-count group means / P95 (Sec.~6.4) & Live run & 32 (8 per peer count) \\
4-vs-32-peer paired comparison (Sec.~6.1.2) & Live run, paired by counterbalancing block & 8 pairs \\
Throughput condition/interaction model (Sec.~6.3) & Live run & 100 (10 levels $\times$ 2 conditions $\times$ 5) \\
Energy single-channel vs.\ residue (Sec.~6.7) & Bench sample & 150 per mode \\
Cryptographic timing (Sec.~6.7) & Bench sample & 300 per operation \\
Field authorization decisions (Sec.~6.1--6.2) & Recorded decision & \AuthDecisions{} \\
Authorization-boundary campaign (Sec.~6.6) & Scripted attempt & \SecurityAttemptsTotal{} \\
\bottomrule
\end{tabular}
\end{table}

\section*{S5. Grant rate by role-operation cell}

\begin{longtable}{llr}
\caption{Grant rate by role-operation cell (all cells in the trace).}\label{tab:s5}\\
\toprule
Role & Operation & Grant rate (\%) \\
\midrule
\endfirsthead
\toprule
Role & Operation & Grant rate (\%) \\
\midrule
\endhead
\bottomrule
\endfoot
Admin & ControlZone & 91.39 \\
Admin & ReadAudit & 93.06 \\
Admin & ReadEnvironmentalRecord & 91.11 \\
Admin & ReadProvenance & 92.06 \\
Admin & ReadSensorData & 93.12 \\
Admin & ReadZone & 92.67 \\
Admin & WriteSensor & 92.73 \\
Agronomist & ControlZone & 92.57 \\
Agronomist & ReadAudit & 91.11 \\
Agronomist & ReadEnvironmentalRecord & 92.45 \\
Agronomist & ReadProvenance & 92.06 \\
Agronomist & ReadSensorData & 92.11 \\
Agronomist & ReadZone & 92.85 \\
Agronomist & WriteSensor & 92.57 \\
Certifier & ControlZone & 91.00 \\
Certifier & ReadAudit & 91.95 \\
Certifier & ReadEnvironmentalRecord & 92.62 \\
Certifier & ReadProvenance & 92.45 \\
Certifier & ReadSensorData & 92.29 \\
Certifier & ReadZone & 92.73 \\
Certifier & WriteSensor & 91.89 \\
Farmer & ControlZone & 90.61 \\
Farmer & ReadAudit & 92.84 \\
Farmer & ReadEnvironmentalRecord & 92.45 \\
Farmer & ReadProvenance & 91.95 \\
Farmer & ReadSensorData & 92.34 \\
Farmer & ReadZone & 92.68 \\
Farmer & WriteSensor & 92.51 \\
Sensor & WriteSensor & 100.00 \\
SupplyChain & ControlZone & 92.11 \\
SupplyChain & ReadAudit & 92.51 \\
SupplyChain & ReadEnvironmentalRecord & 92.28 \\
SupplyChain & ReadProvenance & 92.68 \\
SupplyChain & ReadSensorData & 92.62 \\
SupplyChain & ReadZone & 91.73 \\
SupplyChain & WriteSensor & 92.17 \\
\end{longtable}

\section*{S6. Throughput by concurrency level}

\begin{tabular}{rrrrr}
\toprule
Concurrency & Baseline TPS & HRBAC TPS & \% lower & Holm-adj.\ $p$ \\
\midrule
10 & 42.4 & 37.9 & 10.6 & $<$0.001 \\
20 & 51.2 & 46.0 & 10.3 & $<$0.001 \\
30 & 61.2 & 55.2 & 9.7 & $<$0.001 \\
40 & 70.5 & 63.0 & 10.7 & $<$0.001 \\
50 & 68.2 & 62.2 & 8.9 & $<$0.001 \\
60 & 66.1 & 60.7 & 8.2 & $<$0.001 \\
70 & 63.6 & 59.3 & 6.9 & $<$0.001 \\
80 & 63.4 & 57.4 & 9.4 & $<$0.001 \\
90 & 61.1 & 55.1 & 9.8 & $<$0.001 \\
100 & 59.5 & 54.1 & 9.1 & $<$0.001 \\
\bottomrule
\end{tabular}

\section*{S7. All 32 peer-scaling run-level observations}

\begin{longtable}{rrrrrrr}
\caption{All 32 peer-scaling run-level observations.}\label{tab:s7}\\
\toprule
Serial pos. & Run ID & Peers & Block (rep.) & Total lat.\ (ms) & Span (ms) & P95 (ms) \\
\midrule
\endfirsthead
\toprule
Serial pos. & Run ID & Peers & Block (rep.) & Total lat.\ (ms) & Span (ms) & P95 (ms) \\
\midrule
\endhead
\bottomrule
\endfoot
1 & PEER-04-R1 & 4 & 1 & 1200.4 & 321.0 & 1276.5 \\
2 & PEER-08-R1 & 8 & 1 & 1141.5 & 319.8 & 1239.4 \\
3 & PEER-16-R1 & 16 & 1 & 1057.8 & 323.1 & 1148.7 \\
4 & PEER-32-R1 & 32 & 1 & 810.9 & 316.7 & 889.2 \\
5 & PEER-08-R2 & 8 & 2 & 1146.4 & 319.6 & 1236.0 \\
6 & PEER-16-R2 & 16 & 2 & 1062.2 & 322.5 & 1152.1 \\
7 & PEER-32-R2 & 32 & 2 & 821.7 & 316.5 & 896.0 \\
8 & PEER-04-R2 & 4 & 2 & 1216.1 & 320.9 & 1310.3 \\
9 & PEER-16-R3 & 16 & 3 & 1061.9 & 323.4 & 1148.8 \\
10 & PEER-32-R3 & 32 & 3 & 813.5 & 317.3 & 898.2 \\
11 & PEER-04-R3 & 4 & 3 & 1201.5 & 321.3 & 1289.1 \\
12 & PEER-08-R3 & 8 & 3 & 1154.8 & 320.7 & 1251.2 \\
13 & PEER-32-R4 & 32 & 4 & 812.2 & 317.7 & 910.1 \\
14 & PEER-04-R4 & 4 & 4 & 1200.8 & 321.4 & 1289.2 \\
15 & PEER-08-R4 & 8 & 4 & 1166.7 & 320.4 & 1262.5 \\
16 & PEER-16-R4 & 16 & 4 & 1054.3 & 322.6 & 1149.6 \\
17 & PEER-04-R5 & 4 & 5 & 1201.3 & 321.5 & 1291.1 \\
18 & PEER-32-R5 & 32 & 5 & 808.4 & 317.0 & 888.9 \\
19 & PEER-16-R5 & 16 & 5 & 1062.9 & 323.5 & 1153.8 \\
20 & PEER-08-R5 & 8 & 5 & 1145.5 & 320.7 & 1228.3 \\
21 & PEER-32-R6 & 32 & 6 & 813.5 & 317.4 & 912.3 \\
22 & PEER-16-R6 & 16 & 6 & 1060.2 & 322.1 & 1147.2 \\
23 & PEER-08-R6 & 8 & 6 & 1142.2 & 320.6 & 1238.4 \\
24 & PEER-04-R6 & 4 & 6 & 1199.8 & 321.3 & 1283.9 \\
25 & PEER-16-R7 & 16 & 7 & 1054.9 & 323.2 & 1142.4 \\
26 & PEER-08-R7 & 8 & 7 & 1142.1 & 320.5 & 1225.1 \\
27 & PEER-04-R7 & 4 & 7 & 1203.9 & 321.5 & 1306.6 \\
28 & PEER-32-R7 & 32 & 7 & 809.2 & 316.7 & 898.9 \\
29 & PEER-08-R8 & 8 & 8 & 1158.9 & 321.1 & 1235.1 \\
30 & PEER-04-R8 & 4 & 8 & 1207.1 & 321.6 & 1307.1 \\
31 & PEER-32-R8 & 32 & 8 & 813.3 & 316.8 & 896.8 \\
32 & PEER-16-R8 & 16 & 8 & 1047.5 & 322.4 & 1143.6 \\
\end{longtable}

\section*{S8. Counterbalancing schedule}

Cross-tabulation confirming each configured peer count occupies each within-cycle position exactly twice across the two complementary Latin-square cycles (cycle = repetitions 1--4 vs.\ 5--8).

\begin{center}
\begin{tabular}{rrrr}
\toprule
Peer count & Position in cycle & Cycle & Runs \\
\midrule
4 & 1 & 1 & 1 \\
8 & 1 & 1 & 1 \\
16 & 1 & 1 & 1 \\
32 & 1 & 1 & 1 \\
4 & 2 & 1 & 1 \\
8 & 2 & 1 & 1 \\
16 & 2 & 1 & 1 \\
32 & 2 & 1 & 1 \\
4 & 3 & 1 & 1 \\
8 & 3 & 1 & 1 \\
16 & 3 & 1 & 1 \\
32 & 3 & 1 & 1 \\
4 & 4 & 1 & 1 \\
8 & 4 & 1 & 1 \\
16 & 4 & 1 & 1 \\
32 & 4 & 1 & 1 \\
4 & 1 & 2 & 1 \\
8 & 1 & 2 & 1 \\
16 & 1 & 2 & 1 \\
32 & 1 & 2 & 1 \\
4 & 2 & 2 & 1 \\
8 & 2 & 2 & 1 \\
16 & 2 & 2 & 1 \\
32 & 2 & 2 & 1 \\
4 & 3 & 2 & 1 \\
8 & 3 & 2 & 1 \\
16 & 3 & 2 & 1 \\
32 & 3 & 2 & 1 \\
4 & 4 & 2 & 1 \\
8 & 4 & 2 & 1 \\
16 & 4 & 2 & 1 \\
32 & 4 & 2 & 1 \\
\bottomrule
\end{tabular}
\end{center}

\section*{S9. Revocation stage timing -- withdrawn from V10}

An earlier internal draft reported specific stage-duration statistics (CRL publication mean 60.4\,s, gossip propagation mean 4.21\,min). \texttt{raw/field/revocation\_stages.csv} records a single \texttt{stage\_timestamp} per stage row with no paired start/end timestamp and no reliable per-episode chaining (Section~6.9.4 of the main text), so a stage \emph{duration} cannot be recomputed from the released columns. These specific duration figures are withdrawn in V10 rather than carried forward unverified; Table~4 of the main text and \texttt{analysis/DATA\_DISCREPANCIES.md} record this explicitly. Stage \emph{counts} and the delayed-status proportion (Section~6.5) are unaffected, since they do not require a duration field.

\section*{S10. CRT recovery bounds by missing residue}

\begin{tabular}{lrr}
\toprule
Missing residue pair & Deployed bound & Corrected bound \\
\midrule
$\{97,101\}$ (missing mod 103) & 9,797 & 64,262 \\
$\{97,103\}$ (missing mod 101) & 9,991 & 64,515 \\
$\{101,103\}$ (missing mod 97) & 10,403 & 64,770 \\
\bottomrule
\end{tabular}

\section*{S11. Authorization-boundary denial-mechanism breakdown by scripted scenario}

\begingroup
\small
\resizebox{\linewidth}{!}{%
\begin{tabular}{lrrrrrr}
\toprule
Scenario & CERTIFICATE\_INVALID & OP\_NOT\_PERMITTED & REPLAY\_DETECTED & REVOKED\_CREDENTIAL & ROLE\_INSUFFICIENT & ZONE\_MISMATCH \\
\midrule
brute\_force\_role\_probe & 173 & 164 & 171 & 180 & 136 & 176 \\
expired\_cert\_reuse & 149 & 174 & 165 & 181 & 165 & 166 \\
replay\_attack & 178 & 157 & 162 & 189 & 149 & 165 \\
revoked\_cert\_reuse & 171 & 172 & 177 & 152 & 165 & 163 \\
role\_escalation & 164 & 180 & 143 & 160 & 178 & 175 \\
token\_tampering & 160 & 192 & 184 & 155 & 156 & 153 \\
unauthorized\_resource\_access & 171 & 166 & 164 & 166 & 176 & 157 \\
zone\_boundary\_violation & 165 & 168 & 166 & 157 & 168 & 176 \\
\bottomrule
\end{tabular}%
}
\endgroup

\section*{S12. Data dictionary pointer}

A field-level data dictionary for every released raw trace, including which fields are withdrawn or reinterpreted following the implementation audit, is maintained in \texttt{paper\_iot/V10/analysis/VARIABLE\_MAPPING.md} and \texttt{paper\_iot/V10/analysis/DATA\_PROVENANCE\_REPORT.md}, generated alongside this supplement.

\section*{S13. Evidence needed to reduce each limitation}

\begin{longtable}{p{2.1cm}p{11cm}}
\caption{Evidence that would reduce each principal limitation (Table~5 of the main text).}\label{tab:s13}\\
\toprule
Limitation & What would reduce it \\
\midrule
\endfirsthead
\toprule
Limitation & What would reduce it \\
\midrule
\endhead
\bottomrule
\endfoot
L1 & Instrument one transaction timestamped at endorsement, ordering, validation and commit boundaries, and document the \texttt{rbac\_overhead\_ms} timer's exact start/stop points in source. \\
L2 & Re-run the peer-multiplicity campaign across genuinely distinct physical hosts (not oversubscribed logical peers on four fixed machines), and add a throughput-vs-peer-count arm. \\
L3 & Test additional, randomized concurrency levels (not a single fixed ascending sequence) and instrument CPU, memory, disk and CouchDB service time per component to localise the saturating resource. \\
L4 & Release a permitted-request (positive) corpus alongside the denial corpus, and encode zone, temporal, replay and identity attributes explicitly so the independent oracle can score every dimension, not only role/operation. \\
L5 & Add an episode/correlation identifier carried through every revocation lifecycle stage, and record which post-revocation attempts were made so a success rate is estimable. \\
L6 & Implement, tag and release the corrected chaincode, firmware and gateway verifier described in Section~4.5, then run the full performance and security campaigns against it under the same protocol as the historical measurements. \\
\end{longtable}

\bibliographystyle{plainnat}
\bibliography{../manuscript/references}

\end{document}
